% concatenated from letter.tex
% IEEE Control Systems Letters (L-CSS) version of the RSSB paper.
% Format: two-column IEEE conference format, 6 pages strict (incl. references).
% ieeeconf.cls is a thin wrapper around IEEEtran conference mode; if preferred,
% swap the next line for \documentclass[letterpaper,10pt,conference]{ieeeconf}.
\documentclass[letterpaper,10pt,conference]{IEEEtran}
\usepackage{amsmath,amssymb,amsthm}
\usepackage{graphicx}
\usepackage{hyperref}

\newtheorem{theorem}{Theorem}
\newtheorem{lemma}{Lemma}
\newtheorem{proposition}{Proposition}
\newtheorem{corollary}{Corollary}
\newtheorem{definition}{Definition}
\newtheorem{remark}{Remark}
\newtheorem{example}{Example}

\newcommand{\KL}{\operatorname{KL}}
\newcommand{\EE}{\mathbb{E}}
\newcommand{\Var}{\operatorname{Var}}
\newcommand{\Cov}{\operatorname{Cov}}
\newcommand{\RR}{\mathbb{R}}
\newcommand{\cN}{\mathcal{N}}
\newcommand{\cJ}{\mathcal{J}}



\begin{document}
	\title{The Risk-Sensitive Schr\"odinger Bridge\\ Is Not a KL Projection}
	\author{Hamidreza Behjoo  \IEEEmembership{}
		\thanks{H.~Behjoo is with the Institute of Systems Medicine,
			Chinese Academy of Medical Sciences, Suzhou 215123, China
			(e-mail: hamidreza.behjoo@gmail.com).} }
	
	\maketitle
	\begin{abstract}
		The Schr\"odinger bridge owes its computational power to a single structural
		fact: by Girsanov's theorem the controlled problem is a Kullback--Leibler (KL)
		projection onto a fixed reference measure, solvable by alternating projections.
		This letter shows that the fact does not survive risk sensitivity. When the
		expected path cost is replaced by the entropic risk measure and both endpoint
		marginals are kept as hard constraints, the resulting fixed-point
		bridge value $J_\theta$ (the soft-problem value at the multiplier that enforces the terminal constraint) admits no representation as a constrained KL minimum against any fixed path-space
		reference with a regular endpoint law (a class strictly larger than the
		uniformly elliptic diffusion references: no Markov property is required),
		even allowing an additive
		normalisation depending on the initial marginal. Moreover, no single
		reference generates the one-parameter family in the risk parameter. The obstruction is
		computed in closed form: the Gaussian bridge value violates, by exactly
		$\theta/2$, a heat equation that any Gaussian smoothing of a fixed endpoint
		density must obey. In place of the projection, the theory rests on a
		terminal-multiplier fixed point and an asymmetric factorisation penalising the score energy
		of the backward factor.
	\end{abstract}
	
	\begin{IEEEkeywords}
		Schr\"odinger bridge, risk-sensitive control, entropic risk measure, KL
		duality.
	\end{IEEEkeywords}
	
	\section{Introduction}
	\label{sec:intro}
	
	The Schr\"odinger bridge problem lies at the intersection of stochastic
	control, optimal transport, and generative
	modelling~\cite{Leonard2014,chen2021liaisons}. Its
	computational power rests on a single structural fact: by Girsanov's theorem
	the controlled problem is a Kullback--Leibler (KL) projection onto a fixed
	reference measure, and can therefore be solved by alternating projections or
	Sinkhorn iteration in the static formulation. This letter shows that the projection
	breaks when the expected path cost is replaced by the entropic risk measure
	$\frac1\theta\log\EE\,e^{\theta\,\mathrm{cost}}$, $\theta\in(0,1)$, with
	both endpoint marginals kept as hard constraints. The proof computes the
	obstruction explicitly: the Gaussian bridge value violates, by exactly
	$\theta/2$, a heat equation that the endpoint law of any regular reference
	must obey (Section~\ref{sec:no-kl}).
	
	Risk sensitivity alone would destroy nothing: for a fixed path cost
	the entropic objective differs from a linear functional of the path measure
	only by a monotone transformation, so the constrained problem remains a linear
	program with the usual Schr\"odinger multiplier structure. The obstruction
	here is that the kinetic cost is a functional of the controlled law itself rather than a fixed path functional; jointly with the hard endpoint constraints, this destroys the KL-projection structure on which the classical
	and computational theory rests (Section~\ref{sec:no-kl}). Two values must be
	kept distinct throughout: the {fixed-point bridge value} $J_\theta$ of
	\eqref{eq:bridge-value}, developed here, and the {hard-constrained
		risk-sensitive value} $\cJ_\theta$ of \eqref{eq:hard-rs}, a different object
	for $\theta>0$. The two coincide at $\theta=0$ and separate at first order by
	an explicit variance--covariance gap \eqref{eq:first-order-gap}.
	
	\paragraph*{Contributions}
	(i) a formulation of the risk-sensitive bridge in which the hard terminal
	constraint is enforced by a terminal-multiplier fixed point, with a
	gauge-invariant bridge value $J_\theta$ (Section~\ref{sec:rsb});
	(ii) a three-step impossibility proof (the per-$\theta$ KL representation is circular, no reference is uniform in $\theta$, and no fixed endpoint-regular reference represents $J_\theta$ at all), with the obstruction computed
	explicitly as a $\theta/2$ heat-equation defect (Section~\ref{sec:no-kl});
	(iii) a closed-form Riccati reduction for Gaussian targets to decoupled scalar
	shooting equations, proving well-posedness of the fixed point in that setting
	(Section~\ref{sec:gaussian}); and
	(iv) the structure that replaces the projection: an asymmetric factorisation
	$p^*=\varphi\,\psi^{1/(1-\theta)}$ penalising the score energy of the backward
	factor (Section~\ref{sec:replacement}).
	
	\paragraph*{Related work}
	Exponential-of-integral cost criteria go back to~\cite{jacobson1973} and were
	developed in generality by~\cite{whittle1990,flemingmceneaney1995};
	the equivalence with a minimax game against an entropy-penalised adversary is
	central to robust control~\cite{basar1995}, and
	risk-sensitive extensions within linearly solvable and path-integral control
	appear in~\cite{vandenbroek2010}. All of these consider
	unconstrained problems: the terminal law is free, and the two hard
	marginal constraints responsible for our obstruction are absent. Pricing the
	log-likelihood rather than the kinetic energy~\cite{ito2024renyi} yields a
	R\'enyi divergence: a convex program, but again with a free terminal law. Existing ``robust'' bridge variants relax the hard terminal
	constraint~\cite{chen2019relaxed} or address parameter perturbations; none
	modifies the path-cost objective, and all retain the KL-projection structure. To the best of our
	knowledge this is the first formulation of the Schr\"odinger bridge under the
	entropic risk measure with hard endpoint constraints, and the first
	characterisation of the resulting breakdown of KL/I-projection duality.
	
	\section{Preliminaries: the standard bridge}
	\label{sec:prelim}
	
	Throughout, $\alpha:=1-\theta$, $\cN(m,\Sigma)$ denotes the Gaussian law, and
	$\langle f,\mu\rangle:=\int f\,d\mu$. Fix the horizon $[0,1]$ and consider the controlled diffusion in $\RR^d$
	\begin{equation}
		\label{eq:sde}
		dX_t = \bigl(f_t(X_t)+u_t(X_t)\bigr)\,dt+\sqrt{\kappa_t}\,dW_t,\qquad X_0\sim\mu_0,
	\end{equation}
	with given drift $f_t$, scalar $\kappa_t>0$, Markov control $u$, and $P^u$ the
	path law on $C([0,1];\RR^d)$; $P^0$ is the uncontrolled reference. The
	standard Schr\"odinger bridge~\cite{Leonard2014} is
	\begin{equation}
		\label{eq:sb}
		\inf_{u}\;\EE^{P^u}\!\left[\int_0^1\!\left(\frac{1}{2\kappa_t}\|u_t(X_t)\|^2+V_t(X_t)\right)dt\right]
		\quad\text{s.t.}\;X_1\sim\mu_1.
	\end{equation}
	By Girsanov's theorem the objective equals $\KL(P^u\|P^V)$ up to an additive
	constant, where $dP^V/dP^0\propto\exp(-\int_0^1 V_t(X_t)\,dt)$; the bridge is
	therefore the static problem
	$\inf_{Q:\,Q_0=\mu_0,\,Q_1=\mu_1}\KL(Q\|P^V)$, an I-projection onto a fixed
	reference. The value function $S_0$ of the soft problem associated with \eqref{eq:sb}, the $\theta=0$ member of the risk-sensitive family introduced in Section~\ref{sec:rsb}, satisfies
	\begin{equation}
		\label{eq:hjb-std}
		\begin{split}
			-\partial_t S_0 &= V_t+f_t\cdot\nabla S_0+\frac{\kappa_t}{2}\Delta S_0-\frac{\kappa_t}{2}\|\nabla S_0\|^2,\\
			S_0(1,x)&=\lambda(x),
		\end{split}
	\end{equation}
	with optimal feedback $u^*=-\kappa_t\nabla S_0$, and the Cole--Hopf transform
	$S_0=-\log\psi$ linearises \eqref{eq:hjb-std} to
	\begin{equation}
		\label{eq:psi-bwd}
		-\partial_t\psi+V\psi = f_t\cdot\nabla\psi+\frac{\kappa_t}{2}\Delta\psi,
		\qquad \psi(1,x)=e^{-\lambda(x)}.
	\end{equation}
	With the forward factor $\varphi$ solving
	\begin{equation}
		\label{eq:phi-fwd}
		\partial_t\varphi = -\nabla\cdot(f_t\varphi)+\frac{\kappa_t}{2}\Delta\varphi-V\varphi,
		\qquad \varphi(0,\cdot)=\frac{\mu_0}{\psi(0,\cdot)},
	\end{equation}
	the optimally controlled density admits the symmetric factorisation
	$p^*=\varphi\psi$. We work formally with classical positive solutions of the
	forward--backward system under standard parabolic regularity; the Gaussian
	case of Section~\ref{sec:gaussian} is fully self-contained.
	
	\section{The risk-sensitive bridge}
	\label{sec:rsb}
	
	Replace the linear expectation in \eqref{eq:sb} by the cumulant-generating
	function. For a prescribed terminal cost $\lambda$, the {soft}
	risk-sensitive value is
	\begin{equation}
		\label{eq:rs-obj}
		\begin{split}
			S_\theta(\lambda):=\inf_{u}\frac{1}{\theta}
			\log\EE^{P^u}\!\Bigl[\exp\Bigl(&\theta\!\int_0^1\!\Bigl(\tfrac{\|u_t\|^2}{2\kappa_t}+V_t(X_t)\Bigr)dt\\
			&+\theta\lambda(X_1)\Bigr)\Bigr],
		\end{split}
	\end{equation}
	with only the initial constraint imposed; at $\theta=0$ this is the classical
	soft value, and with the terminal constraint restored the bridge value
	\eqref{eq:bridge-value} recovers \eqref{eq:sb} exactly. The {bridge} problem imposes in addition the hard terminal
	constraint $X_1\sim\mu_1$, and $\lambda$ is then no longer a datum: it becomes a
	{terminal multiplier}, i.e.\ a shooting parameter, enforcing the constraint.
	
	\subsection{Gauge and the bridge value}
	The soft value is exactly shift-equivariant,
	$S_\theta(\lambda+c)=S_\theta(\lambda)+c$ for $c\in\RR$, while the terminal coupling
	\eqref{eq:Psi1} below is blind to constants; the multiplier is determined only
	up to the coset $\lambda^*+\RR$, and the optimally controlled process is
	untouched by the constant. The gauge-invariant combination is the difference
	\begin{equation}
		\label{eq:bridge-value}
		J_\theta(\mu_0\to\mu_1):=S_\theta(\lambda^*)-\langle\lambda^*,\mu_1\rangle,
	\end{equation}
	with $\lambda^*$ any representative of the coset.
	
	\begin{definition}[Bridge value]
		\label{def:bridge-value}
		The {risk-sensitive bridge value} is the gauge-invariant fixed-point
		value \eqref{eq:bridge-value}. For computations we fix the representative by
		the gauge $\langle\lambda^*,\mu_1\rangle=0$, under which
		$J_\theta=S_\theta(\lambda^*)$; the definition does not depend on the gauge.
	\end{definition}
	
	The definition is canonical in two senses. At $\theta=0$ the terminal
	coupling becomes the classical Schr\"odinger system and
	\eqref{eq:bridge-value} is precisely the Legendre identity
	$J_0=\sup_{\lambda}\{S_0(\lambda)-\langle\lambda,\mu_1\rangle\}$: the classical bridge
	value is recovered with no arbitrary choices. Second, the fixed point is
	well-posed at least in the Gaussian setting, by strict monotonicity of the
	scalar shooting map (Remark~\ref{rem:shoot-wp}). For $\theta>0$ the Legendre
	identity fails: the first variation of
	$\lambda\mapsto S_\theta(\lambda)-\langle\lambda,\mu_1\rangle$ at the fixed point is
	$q_1^\theta-\mu_1$, with $q_1^\theta$ a {tilted} terminal law, so the fixed point is not a saddle point of a Lagrangian.
	
	\subsection{The hard value and the first-order gap}
	For $\theta>0$, \eqref{eq:bridge-value} is not the value of the
	hard-constrained risk-sensitive control problem
	\begin{equation}
		\label{eq:hard-rs}
		\begin{split}
			\cJ_\theta(\mu_0\to\mu_1)
			:=\inf_{u:\,X_0\sim\mu_0,\,X_1\sim\mu_1}&\frac1\theta\log
			\EE^{u}\!\bigl[e^{\theta K^u}\bigr],\\
			K^u&=\int_0^1\Bigl(\tfrac{\|u_t\|^2}{2\kappa_t}+V_t\Bigr)dt .
		\end{split}
	\end{equation}
	The two coincide at $\theta=0$ and separate at first order: writing $u_0^*$
	for the classical bridge control and $\lambda_0^*$ for the classical multiplier,
	a cumulant expansion with an envelope argument gives
	\begin{equation}
		\label{eq:first-order-gap}
		J_\theta-\cJ_\theta
		=\theta\Bigl(\tfrac12\Var^{u_0^*}\!\bigl(\lambda_0^*(X_1)\bigr)
		+\Cov^{u_0^*}\!\bigl(K,\,\lambda_0^*(X_1)\bigr)\Bigr)+o(\theta),
	\end{equation}
	generically nonzero, of either sign. The same expansion applied to
	$J_\theta$ alone gives
	\begin{equation}
		\label{eq:first-order-value}
		J_\theta=J_0+\frac{\theta}{2}\,
		\Var^{u_0^*}\!\bigl(K+\lambda_0^*(X_1)\bigr)+o(\theta),
	\end{equation}
	and \eqref{eq:first-order-gap} is its difference with the companion expansion
	of $\cJ_\theta$. The two problems differ because the kinetic cost is a
	functional of the controlled law, not a fixed path functional: in
	\eqref{eq:hard-rs} the constraint is on the law of $X_1$, so its KKT
	multiplier is not absorbed into the path cost and acts outside the entropic
	exponential, and \eqref{eq:hard-rs} is not a linear program over path
	measures. The fixed-point construction \eqref{eq:bridge-value} instead keeps
	the terminal potential inside the exponential, which is precisely what makes
	the generalised Cole--Hopf transform below available.
	
	\subsection{HJB equation and generalised Cole--Hopf}
	By dynamic programming, the state-conditioned soft value $S_\theta(t,x)$
	satisfies
	\begin{equation}
		\label{eq:hjb-rs}
		\begin{split}
			-\partial_t S_\theta &= V_t+f_t\cdot\nabla S_\theta+\frac{\kappa_t}{2}\Delta S_\theta
			-\frac{\kappa_t}{2}(1-\theta)\|\nabla S_\theta\|^2,\\
			S_\theta(1,x)&=\lambda(x).
		\end{split}
	\end{equation}
	The generalised Cole--Hopf transform
	\begin{equation}
		\label{eq:ch-rs}
		\psi_t(x):=\exp\!\bigl(-(1-\theta)S_\theta(t,x)\bigr)
		\;\Longleftrightarrow\;
		S_\theta(t,x)=-\frac{\log\psi_t(x)}{1-\theta}
	\end{equation}
	linearises \eqref{eq:hjb-rs} to the backward PDE
	\begin{equation}
		\label{eq:psi-rs}
		\begin{split}
			-\partial_t\psi_t+(1-\theta)V_t\psi_t
			&= f_t\cdot\nabla\psi_t+\frac{\kappa_t}{2}\Delta\psi_t,\\
			\psi(1,x)&=e^{-(1-\theta)\lambda(x)},
		\end{split}
	\end{equation}
	whose generator is $\theta$-independent. The optimal feedback is
	\begin{equation}
		\label{eq:u-star}
		u_t^*(x)=\frac{\kappa_t}{1-\theta}\nabla\log\psi_t(x).
	\end{equation}
	
	\subsection{Asymmetric factorisation}
	Under \eqref{eq:u-star}, the controlled density $p^*$ satisfies the
	Fokker--Planck equation with drift
	$f_t+\frac{\kappa_t}{1-\theta}\nabla\log\psi_t$.
	
	\begin{proposition}[Forward decoupling]
		\label{prop:decouple}
		Let $\psi$ solve \eqref{eq:psi-rs} and let $p^*$ be the optimally controlled
		density. Then
		\begin{equation}
			\label{eq:factor-rs}
			p^*(t,x)=\varphi(t,x)\,\psi(t,x)^{\frac{1}{1-\theta}},
		\end{equation}
		where $\varphi$ satisfies the closed linear forward equation
		\begin{equation}
			\label{eq:G-fwd}
			\partial_t \varphi = -\nabla\cdot(f_t \varphi)+\frac{\kappa_t}{2}\Delta \varphi-V_{\rm eff}(t,x)\varphi,
		\end{equation}
		with effective potential
		\begin{equation}
			\label{eq:Veff}
			V_{\rm eff}(t,x)=V(t,x)+\frac{\kappa_t\theta}{2(1-\theta)^2}\|\nabla\log\psi(t,x)\|^2,
		\end{equation}
		and initial condition $\varphi(0,x)=\mu_0(x)/\psi(0,x)^{1/(1-\theta)}$.
	\end{proposition}
	
	\begin{proof}
		Insert the ansatz $p^*=\varphi\psi^r$ into the Fokker--Planck equation. The
		drift divergence produces a cross term proportional to
		$\nabla\varphi\cdot\nabla\psi$ with coefficient $r-(1-\theta)^{-1}$; setting
		it to zero yields $r=1/(1-\theta)$. Dividing by $\psi^r$ and using
		\eqref{eq:psi-rs} to eliminate $\partial_t\psi$ and $\Delta\psi$, the
		remaining scalar terms combine to give \eqref{eq:G-fwd}--\eqref{eq:Veff}.
	\end{proof}
	
	Since $\theta\ge0$, \eqref{eq:Veff} gives $V_{\rm eff}\ge V$ pointwise,
	strictly wherever $\nabla\log\psi\neq0$: risk sensitivity can only add to the
	potential landscape seen by the forward factor, concentrated where the
	backward score is steep. The hard terminal constraint $p^*(1,\cdot)=\mu_1$
	reads
	\begin{equation}
		\label{eq:Psi1}
		\varphi(1,x)\,\psi(1,x)^{1/(1-\theta)}=\mu_1(x),
	\end{equation}
	equivalently $\varphi(1,\cdot)\,e^{-\lambda^*(\cdot)}=\mu_1$ in terms of the
	multiplier: the compact identity that makes $\lambda^*$ a multiplier for the hard
	constraint, determined implicitly as a shooting parameter. At $\theta=0$ the
	score-energy term vanishes and \eqref{eq:G-fwd}--\eqref{eq:Psi1} reduce to
	\eqref{eq:phi-fwd} and the symmetric factorisation.
	
	\section{The fixed-point value has no KL dual}
	\label{sec:no-kl}
	
	For the standard bridge, Girsanov yields
	$\EE^{P^u}\!\int_0^1(\frac{\|u_t\|^2}{2\kappa_t}+V_t)\,dt
	=\KL(P^u\|P^V)+\text{const}$, so the bridge is the I-projection
	$\inf_{Q:\,Q_0=\mu_0,\,Q_1=\mu_1}\KL(Q\|P^V)$. It is natural to ask whether
	$J_\theta$, for $\theta>0$, admits an analogous projection representation. It
	does not, in a precise sense. The statements concern the fixed-point value
	$J_\theta$; the hard value $\cJ_\theta$ likewise admits no linear-programming
	formulation, for the independent reason that its exponent is law-dependent. The structural obstruction is the asymmetric
	factorisation \eqref{eq:factor-rs}: the Gibbs form
	$dQ^*/dR\propto\varphi\cdot\psi$, with factors of equal weight, is exactly
	what underlies the static KL dual, and it is what $\theta>0$ breaks. The asymmetry alone is not a proof: Lemma~\ref{lem:match} shows that a
	solution-dependent reference reproduces the risk-sensitive bridge optimizer in
	full; the
	impossibility is argued at the level of the value.
	
	\begin{theorem}[No KL dual]
		\label{thm:flagship}
		Fix the problem data $(f,\kappa,V)$, let $\theta\in(0,1)$, and assume the
		terminal coupling \eqref{eq:Psi1} admits a solution $\lambda^*+\RR$ for the
		endpoint pairs under consideration, so that $J_\theta$ is defined; this
		holds in the Gaussian setting, which is what the proofs test against.
		Then:
		(i) {(circularity)} for each fixed $\theta$, a Markov reference whose
		I-projection value reproduces $J_\theta$ exists only if the reference is
		built from the unknown solution $\nabla S_\theta$ itself, via the matching
		condition \eqref{eq:match};
		(ii) {(no $\theta$-uniform reference)} no single Markov reference
		represents $J_\theta$ as a constrained KL minimum for two distinct values
		of $\theta$;
		(iii) {(no fixed reference at all)} there is no endpoint-regular
		(two-sided Gaussian endpoint joint law, Definition~\ref{def:endpoint-regular})
		path-space reference $R_\theta$, Markovian or not, and no additive
		correction $C(\mu_0)$, such that
		$J_\theta(\mu_0\to\mu_1)=\inf_{Q:\,Q_0=\mu_0,\,Q_1=\mu_1}\KL(Q\|R_\theta)+C(\mu_0)$
		for all absolutely continuous marginals with full support and finite second
		moment. Part (i) is Lemma~\ref{lem:match}, part (ii) is
		Proposition~\ref{prop:no-uniform}, and part (iii) is
		Proposition~\ref{prop:no-kl}; the last two are proved for the scalar
		Brownian data $f\equiv0$, $\kappa=1$, $V\equiv0$ of
		Lemma~\ref{lem:value-pm}: the obstruction appears already in the simplest
		setting, which precludes any representation serving general data.
	\end{theorem}
	
	\subsection{The per-$\theta$ representation is circular}
	Consider a candidate Markov reference $P^{h,U}$ with drift $h$ and
	Feynman--Kac potential $U$ (a sub-probability path measure when $U\neq0$, to
	which $\KL$ extends verbatim); relative to it, $\inf_Q\KL(Q\|P^{h,U})$ is itself
	an ordinary bridge problem whose value function $S_0^{h,U}$ satisfies
	\eqref{eq:hjb-std} with $(f,V)$ replaced by $(h,U)$.
	
	\begin{lemma}[Matching condition]
		\label{lem:match}
		Fix the problem data $(f,\kappa,V)$, a terminal cost $\lambda$, and a Markov
		reference $(h,U)$ for which the two HJB equations admit classical solutions.
		Then $S_0^{h,U}(\cdot\,;\lambda)\equiv S_\theta(\cdot\,;\lambda)$ if and only if
		\begin{equation}
			\label{eq:match}
			\begin{split}
				U(t,x)=V(t,x)&+\bigl(f(t,x)-h(t,x)\bigr)\cdot\nabla S_\theta(t,x)\\
				&+\frac{\kappa_t\theta}{2}\|\nabla S_\theta(t,x)\|^2
			\end{split}
		\end{equation}
		for all $(t,x)$. When \eqref{eq:match} holds with $h=f$, i.e.\
		$U=V_{\rm eff}$ of \eqref{eq:Veff}, and $\lambda=\lambda^*$ implements
		$X_1\sim\mu_1$ via \eqref{eq:Psi1}, the I-projection of $P^{f,V_{\rm eff}}$
		onto $\{Q_0=\mu_0,\,Q_1=\mu_1\}$ coincides with the risk-sensitive bridge optimizer $P^{u^*}$
		itself, not merely with its value.
	\end{lemma}
	
	\begin{proof}
		Subtracting the two HJB equations and using $S_0\equiv S_\theta$ gives
		\eqref{eq:match}; conversely, under \eqref{eq:match} the function
		$S_\theta$ solves the reference HJB with the same terminal datum, and
		uniqueness gives $S_0\equiv S_\theta$. For the last
		statement, the reference bridge's optimal total drift is
		$h-\kappa_t\nabla S_0$; with $h=f$ and $S_0=S_\theta$ this is
		$f+\frac{\kappa_t}{1-\theta}\nabla\log\psi$, exactly \eqref{eq:u-star}.
		Same initial law and same feedback give the same path measure.
	\end{proof}
	
	So a KL-projection representation of $J_\theta$ exists for each fixed $\theta$, but only by building the reference out of $\nabla S_\theta$ itself,
	i.e.\ out of the unknown solution: a construction that presupposes the answer
	and is not a duality in any useful sense; read positively, the matching
	condition made constructive is a fixed-point iteration on the multiplier
	coset.
	
	\subsection{Two ingredients}
	The remaining proofs rest on an explicit Gaussian value and a static
	reduction.
	
	\begin{lemma}[Point-mass Gaussian value]
		\label{lem:value-pm}
		Let $d=1$, $f\equiv0$, $V\equiv0$, $\kappa=1$, $\mu_0=\delta_0$,
		$\mu_1=\cN(m,s^2)$, and write $\alpha=1-\theta$. Then
		\begin{equation}
			\label{eq:value-pm}
			\begin{split}
				v_\theta(m,s^2)&:=J_\theta(\delta_0\to\cN(m,s^2))\\
				&=\frac{1}{2\alpha}\log(1+a_1)-\frac{a_1}{2\alpha}\bigl(m^2+s^2\bigr)\\
				&\qquad+\frac{m^2}{M}-\frac{\alpha B}{2M^2}\,m^2,
			\end{split}
		\end{equation}
		where $a_1$ is determined by the point-mass shooting map
		\begin{equation}
			\label{eq:shoot-pm}
			s^2=\int_0^1\bigl(1+a_1 \tau\bigr)^{-2/\alpha}\,d\tau,
		\end{equation}
		and $M:=\int_0^1(1+a_1\tau)^{-1/\alpha-1}\,d\tau$,
		$B:=\int_0^1(1+a_1\tau)^{-2}\,d\tau=1/(1+a_1)$. The map $a_1\mapsto s^2$ is
		a strictly decreasing bijection of $(-1,\infty)$ onto $(0,\infty)$, with
		$s^2=1\iff a_1=0$. At $\theta=0$, \eqref{eq:value-pm} collapses to
		$\frac12(m^2+s^2-1-\log s^2)=\KL(\cN(m,s^2)\,\|\,\cN(0,1))$, as it must.
	\end{lemma}
	
	\begin{proof}
		Set $k=0$ in the scalar Riccati system of Section~\ref{sec:gaussian}:
		$\dot A=A^2$, $\dot b=Ab$, whose solutions are rational in $t$; passing
		$\sigma(0)\downarrow0$ in the shooting equation \eqref{eq:shoot} yields
		\eqref{eq:shoot-pm} (cf.\ Remark~\ref{rem:shoot-wp}), and direct evaluation of
		\eqref{eq:bridge-value} gives \eqref{eq:value-pm}. Moreover
		$\partial s^2/\partial a_1=-\frac{2}{\alpha}\int_0^1\tau(1+a_1\tau)^{-2/\alpha-1}d\tau<0$
		is nonvanishing, so $a_1$ is a $C^1$ function of $s^2$ down to $s^2=1$,
		with derivative $-\alpha$ there.
	\end{proof}
	
	\begin{definition}[Endpoint regularity]
		\label{def:endpoint-regular}
		A probability measure $R$ on $C([0,1];\RR^d)$ is {endpoint-regular} if
		its endpoint law $R_{01}:=\operatorname{Law}_R(X_0,X_1)$ has a positive
		continuous density $r(x,y)$ with two-sided Gaussian bounds:
		there exist $c,\Lambda>0$ and a continuous $m_0$ of at most linear growth
		such that
		\begin{equation}
			\label{eq:endpoint-regular}
			c^{-1}e^{-\Lambda|y-m_0(x)|^2}\;\le\; r(x,y)\;\le\; c\,e^{-\Lambda^{-1}|y-m_0(x)|^2}
		\end{equation}
		for all $x,y\in\RR^d$. This holds for the path law of any uniformly elliptic diffusion with
		bounded measurable drift and bounded Feynman--Kac potential by Aronson's
		two-sided estimates~\cite{aronson1967}, and for Ornstein--Uhlenbeck
		references by their explicit Gaussian law.
	\end{definition}
	
	\begin{lemma}[Static reduction and the point-mass limit]
		\label{lem:static}
		Let $R$ be endpoint-regular and set
		$I_R(\mu_0,\mu_1):=\inf_{Q:\,Q_0=\mu_0,\,Q_1=\mu_1}\KL(Q\|R)$. Then:
		(i) {(static reduction)}
		$I_R(\mu_0,\mu_1)=\inf_{\pi\in\Pi(\mu_0,\mu_1)}\KL(\pi\|R_{01})$, a static
		Schr\"odinger problem on $\RR^d\times\RR^d$;
		(ii) {(point-mass limit)} if $\mu_0^\varepsilon=\cN(0,\varepsilon I_d)$
		and $\mu_1,\mu_1^*$ are absolutely continuous with finite second moment and
		finite entropy, then, with $\rho_0(\cdot):=R(X_1\in\cdot\mid X_0=0)$,
		\begin{equation}
			\label{eq:pm-limit}
			I_R(\mu_0^\varepsilon,\mu_1)-I_R(\mu_0^\varepsilon,\mu_1^*)
			\;\xrightarrow[\varepsilon\downarrow0]{}\;
			\KL(\mu_1\|\rho_0)-\KL(\mu_1^*\|\rho_0).
		\end{equation}
	\end{lemma}
	
	\begin{proof}
		Part (i) is the chain rule for KL divergence disintegrated over the
		endpoint pair: the bridge term is minimised at $0$ by adopting the bridges
		of $R$. Part (ii) is a two-sided squeeze under
		\eqref{eq:endpoint-regular}: the two-sided Gaussian bounds control the
		contribution of $\{X_0\neq0\}$ as $\varepsilon\downarrow0$, and positivity
		and continuity of $r$ near $(0,\cdot)$ give the matching lower bound.
	\end{proof}
	
	Some restriction on the additive freedom $C$ is necessary for the question
	to be meaningful: if $C$ could depend on both marginals, any reference
	would represent $J_\theta$ by setting
	$C(\mu_0,\mu_1):=J_\theta-\inf\KL(\cdot\,\|R_\theta)$. Dependence on $\mu_0$
	only is the natural class: for the standard bridge the additive term is the
	Feynman--Kac normalisation, independent of the marginals altogether.
	
	\subsection{No reference works uniformly in $\theta$}
	\begin{proposition}[No $\theta$-uniform KL dual]
		\label{prop:no-uniform}
		There is no Markov reference $(h,U)$, chosen independently of $\theta$, with
		$h$ Lipschitz in $x$ of at most linear growth and $U$ bounded, and constants
		$C_\theta$ independent of the marginals, such that, in the setting of
		Lemma~\ref{lem:value-pm},
		$J_\theta(\mu_0\to\mu_1)
		=\inf_{Q:\,Q_0=\mu_0,\,Q_1=\mu_1}\KL(Q\|P^{h,U})+C_\theta$
		holds for two distinct values $\theta\in[0,1)$ and all absolutely continuous
		marginals $\mu_0,\mu_1$.
	\end{proposition}
	
	\begin{proof}
		Suppose a fixed pair $(h,U)$ satisfies the identity for
		$\theta_1\neq\theta_2$; write $R=P^{h,U}$. Apply the identity with
		$\mu_0^\varepsilon=\cN(0,\varepsilon)$ and two Gaussian targets, subtract,
		and let $\varepsilon\downarrow0$. On the left the Riccati formulas of
		Lemma~\ref{lem:value-pm} are continuous in the initial variance down to
		$\sigma(0)=0$; on the right Lemma~\ref{lem:static} applies, since the
		endpoint kernel of $P^{h,U}$ is endpoint-regular (bounded $h$ by Aronson's
		estimate~\cite{aronson1967}; linear $h$ explicitly Gaussian; general
		linear-growth Lipschitz $h$ by Girsanov comparison against the driftless
		kernel and the Cauchy--Schwarz inequality on a finite horizon). Hence, with
		$\rho_0(\cdot):=R(X_1\in\cdot\mid X_0=0)$,
		$v_{\theta_i}(m,s^2)=\KL(\cN(m,s^2)\,\|\,\rho_0)+C_{\theta_i}$ for $i=1,2$
		and all $m\in\RR$, $s^2>0$. Subtracting, $v_{\theta_1}-v_{\theta_2}$ must
		be the constant $C_{\theta_1}-C_{\theta_2}$. But at $s^2=1$ (so $a_1=0$, $M=B=1$),
		\eqref{eq:value-pm} gives
		$v_\theta(m,1)=m^2-\frac{\alpha}{2}m^2=\frac{1+\theta}{2}m^2$, hence
		$v_{\theta_1}(m,1)-v_{\theta_2}(m,1)=\tfrac12(\theta_1-\theta_2)m^2$, which
		is not constant in $m$, a contradiction.
	\end{proof}
	
	Thus no fixed $(h,U)$ generates the family $\{J_\theta\}_{\theta\in[0,1)}$ the
	way the ordinary bridge is a projection onto the single reference $P^V$: the
	risk-sensitive bridge is a one-parameter family of different control problems, not a re-weighting of
	one reference measure.
	
	\subsection{No fixed reference works at all}
	\begin{proposition}[No endpoint-regular reference]
		\label{prop:no-kl}
		Let $d=1$, $f\equiv0$, $\kappa=1$, $V\equiv0$, and let $\theta\in(0,1)$.
		There exists no endpoint-regular probability measure $R_\theta$ on
		$C([0,1];\RR)$ and no map $\mu_0\mapsto C(\mu_0)$ such that
		$J_\theta(\mu_0\to\mu_1)=\inf_{Q:\,Q_0=\mu_0,\,Q_1=\mu_1}\KL(Q\|R_\theta)+C(\mu_0)$
		for all absolutely continuous marginals $\mu_0,\mu_1$ with full support and
		finite second moment.
	\end{proposition}
	
	\begin{proof}
		Suppose $(R_\theta,C)$ satisfies the identity. {Step 1: point-mass
			test.} Apply the identity with $\mu_0^\varepsilon=\cN(0,\varepsilon)$ and
		two Gaussian targets, and subtract: $C$ cancels since it depends only on
		$\mu_0$. Let $\varepsilon\downarrow0$. On the left, the Riccati formulas
		are continuous down to $\sigma(0)=0$; on the right,
		Lemma~\ref{lem:static}(ii) applies with
		$\rho_0(\cdot):=R_\theta(X_1\in\cdot\mid X_0=0)$, giving
		\begin{equation}
			\label{eq:kl-form}
			v_\theta(m,s^2)=\KL\bigl(\cN(m,s^2)\,\|\,\rho_0\bigr)+\mathrm{const}
		\end{equation}
		for all $m\in\RR$, $s^2>0$, with the constant independent of $(m,s^2)$. {Step 2: a deconvolution
			criterion.} Since
		$\KL(\cN(m,s^2)\|\rho)=-\tfrac12\log(2\pi s^2)-\tfrac12
		-\EE_{Y\sim\cN(m,s^2)}[\log\rho(Y)]$, the identity \eqref{eq:kl-form} says
		that
		$L_\theta(m,s^2):=-v_\theta(m,s^2)-\tfrac12\log(2\pi s^2)-\tfrac12$
		is, up to a constant, the Gaussian smoothing at bandwidth $s^2$ of the
		fixed function $\log\rho_0$. Gaussian smoothing obeys the heat
		equation in $(m,s^2)$:
		\begin{equation}
			\label{eq:heat}
			\partial_{s^2}L_\theta=\tfrac12\,\partial_m^2 L_\theta
			\qquad\text{for all }m\in\RR,\ s^2>0,
		\end{equation}
		a necessary condition on the value $v_\theta$ alone. {Step 3: explicit
			failure at $s^2=1$.} At $s^2=1$ one has $a_1=0$, $M=B=1$, so
		$v_\theta(m,1)=\frac{1+\theta}{2}m^2$ and $\tfrac12\partial_m^2
		L_\theta(m,1)=-\tfrac12\partial_m^2 v_\theta(m,1)=-\frac{1+\theta}{2}$.
		On the other hand $v_\theta(0,s^2)=\frac{1}{2\alpha}\log(1+a_1)-\frac{a_1
			s^2}{2\alpha}$, and differentiating \eqref{eq:shoot-pm} at $a_1=0$ gives
		$\frac{da_1}{ds^2}\big|_{s^2=1}=-\alpha$; the two terms of
		$\partial_{s^2}v_\theta(0,s^2)$ then cancel at $s^2=1$, giving
		$\partial_{s^2}L_\theta|_{(0,1)}=-\partial_{s^2}v_\theta|_{(0,1)}
		-\tfrac12=-\frac12$. The heat equation \eqref{eq:heat} would require
		$-\tfrac12=-\tfrac{1+\theta}{2}$, which fails for every $\theta>0$: the
		heat-equation defect is exactly $\theta/2$. Hence no density $\rho_0$, and
		{a fortiori} no reference $R_\theta$, can satisfy \eqref{eq:kl-form}.
		At $\theta=0$ the two sides agree and \eqref{eq:kl-form} holds with
		$\rho_0=\cN(0,1)$, the Brownian endpoint law: the classical KL duality, recovered as a sanity check.
	\end{proof}
	
	\begin{remark}[An inf--sup in place of the KL dual]
		\label{rem:dv}
		The structural reason is visible in the Donsker--Varadhan
		formula~\cite{donsker1975,follmerschied2016}: the entropic risk is itself
		a Legendre object, so $J_\theta$ is an inf--sup of KL-type objects (classically, a game against a drift-perturbing adversary~\cite{jacobson1973,basar1995}), whereas the standard
		bridge is a single inf; a projection would force the inner
		supremum to collapse, which happens only at $\theta=0$.
	\end{remark}
	
	\section{The Gaussian case: closed-form Riccati reduction}
	\label{sec:gaussian}
	
	Specialise to $f_t=0$, $\kappa_t=\kappa>0$, $V(x)=\frac{k}{2}\|x\|^2$
	($k\ge0$), $\mu_0=\cN(0,I_d)$, $\mu_1=\cN(m_1,\Sigma_1)$ with
	$\Sigma_1\in\mathbb{S}_{++}^d$, and define
	$\gamma:=\sqrt{\alpha k/\kappa}$, $\omega:=\sqrt{\alpha k\kappa}$. Posit the
	log-quadratic ansatz
	$\psi(t,x)=\exp(-\frac12 x^\top A(t)x+b(t)^\top x+c(t))$,
	$\varphi(t,x)=\exp(-\frac12 x^\top D(t)x+E(t)^\top x+F(t))$ with
	$A,D\in\mathbb{S}^d$. The optimally controlled marginal is then Gaussian
	with precision $\Sigma(t)^{-1}=D(t)+A(t)/\alpha$; positivity is required only
	of this sum.
	Inserting the ansatz into \eqref{eq:psi-rs} gives the backward Riccati system
	\begin{equation}
		\label{eq:riccati-A}
		\dot A=\kappa A^2-\alpha k I,\qquad
		\dot b=\kappa Ab,
	\end{equation}
	with $\dot c=\frac{\kappa}{2}(\operatorname{tr}A-\|b\|^2)$.
	and since $\|\nabla\log\psi\|^2$ is quadratic, the effective potential
	\eqref{eq:Veff} remains quadratic: the ansatz is self-consistent. The forward
	equation \eqref{eq:G-fwd} yields
	\begin{equation}
		\label{eq:riccati-D}
		\dot D=-\kappa D^2+kI+\frac{\kappa\theta}{\alpha^2}A^2,\qquad
		\dot E=-\kappa DE+\frac{\kappa\theta}{\alpha^2}Ab,
	\end{equation}
	and in the total precision the $kI$ terms cancel:
	\begin{equation}
		\label{eq:riccati-Lambda}
		\frac{d}{dt}\Sigma(t)^{-1}=-\kappa D^2+\frac{\kappa}{\alpha^2}A^2,
	\end{equation}
	a closed system for $(A,\Sigma^{-1})$, which recovers
	the covariance equation
	\begin{equation}
		\label{eq:cov-ode}
		\dot\Sigma=\kappa I-\frac{\kappa}{\alpha}\bigl(A\Sigma+\Sigma A\bigr),
	\end{equation}
	exactly the moment evolution of the controlled Ornstein--Uhlenbeck SDE with
	drift $-\frac{\kappa}{\alpha}(A(t)x-b(t))$.
	
	In the eigenbasis of $\Sigma_1$, $A(t)$ and $\Sigma(t)$ are simultaneously
	diagonal; each diagonal entry satisfies the scalar system
	\begin{equation}
		\label{eq:scalar}
		\dot a_i=\kappa a_i^2-\alpha k,\qquad
		\dot\sigma_i=\kappa-\frac{2\kappa}{\alpha}a_i\sigma_i,
	\end{equation}
	with $\sigma_i(0)=1$, $\sigma_i(1)=\ell_i$. The backward equation has an
	explicit hyperbolic solution, and the terminal condition reduces to the scalar
	shooting equation
	\begin{equation}
		\label{eq:shoot}
		\ell_i=\gamma^{2/\alpha}\left(g_i(1)^{-2/\alpha}
		+\kappa\int_0^1 g_i(\tau)^{-2/\alpha}\,d\tau\right),
	\end{equation}
	with $g_i(\tau)=\gamma\cosh(\omega\tau)+a_{i,1}\sinh(\omega\tau)$ and
	$a_{i,1}:=a_i(1)$: one transcendental equation per coordinate, completely
	decoupled. The mean condition is then a single affine equation per coordinate
	for $b_{i,1}$.
	
	\begin{remark}[Well-posedness; Brownian limit]
		\label{rem:shoot-wp}
		The map $a_{i,1}\mapsto\ell_i$ in \eqref{eq:shoot} is strictly decreasing (for $\tau>0$, $g_i(\tau)$ is strictly increasing in $a_{i,1}$) and surjective onto $(0,\infty)$, so a unique root exists for
		every $\ell_i>0$. This
		proves existence and uniqueness of the terminal-multiplier coset of
		Definition~\ref{def:bridge-value} in the Gaussian setting. In the Brownian
		case $k=0$ the hyperbolic formula degenerates to a rational one, and
		passing $\sigma_i(0)\downarrow0$ recovers the point-mass map
		\eqref{eq:shoot-pm} used in Section~\ref{sec:no-kl}.
	\end{remark}
	
	\paragraph*{The obstruction, in factor language}
	In the same Brownian setting, write the terminal forward factor as
	$\varphi(1,x)\propto\exp(-\tfrac12 D(1)x^2+E(1)x)$. Precision addition in the
	coupling \eqref{eq:Psi1} and the shooting map \eqref{eq:shoot-pm} give
	$D(1)=\frac{1}{s^2}-\frac{a_1}{\alpha}
	=1+\frac{\theta}{3(1-\theta)^2}a_1^2+O(a_1^3)$. For $\theta>0$ the terminal
	precision from a point mass varies with the target variance near $s^2=1$,
	whereas a fixed reference forces target-independent endpoint bridges
	($\theta=0$: $D(1)\equiv1$). This is the factor-level form of the value-level
	defect of Proposition~\ref{prop:no-kl}.
	
	
	\begin{figure}[t]
		\centering
		\includegraphics[width=\columnwidth]{fig_gaussian.pdf}
		\caption{Gaussian bridge with $\kappa=k=1$, $\mu_0=\cN(0,1)$,
			$\mu_1=\cN(1,1/4)$, computed from \eqref{eq:scalar}--\eqref{eq:shoot}.
			(a) Marginal variance $\Sigma(t)$: the endpoint constraints are met
			exactly for every $\theta$, while the intermediate variance contracts as
			$\theta$ grows. (b) Bridge value $J_\theta$: it leaves the classical
			value $J_0$ with the first-order slope of \eqref{eq:first-order-value} and
			diverges at a finite $\theta$, the classical breakdown of the entropic
			criterion~\cite{whittle1990}.}
		\label{fig:gaussian}
	\end{figure}
	
	\paragraph*{Numerical illustration}
	Figure~\ref{fig:gaussian} shows the one-dimensional case $\kappa=k=1$,
	$\mu_0=\cN(0,1)$, $\mu_1=\cN(1,1/4)$, solved through the scalar system
	\eqref{eq:scalar}--\eqref{eq:shoot}. The terminal constraint is met exactly
	at every $\theta$, while the intermediate variance contracts as $\theta$
	grows. The bridge
	value leaves the classical value with the first-order slope of
	\eqref{eq:first-order-value}, then grows superlinearly and diverges at a finite
	$\theta$: for an unbounded initial law the entropic criterion has a finite
	breakdown point, beyond which no finite-cost bridge exists. The breakdown
	attaches to the unbounded initial law rather than to the construction: the
	state-conditioned soft value $S_\theta(t,x)$ and the point-mass case of
	Lemma~\ref{lem:value-pm} remain finite for every $\theta<1$.
	
	\section{What replaces the projection}
	\label{sec:replacement}
	
	The failure of the KL projection is not the failure of structure. For
	$\theta>0$ the fixed-point bridge retains the generalised Cole--Hopf
	linearisation \eqref{eq:ch-rs}--\eqref{eq:psi-rs}, whose backward generator is
	$\theta$-independent, and the asymmetric factorisation
	$p^*=\varphi\,\psi^{1/(1-\theta)}$ under the effective potential
	\eqref{eq:Veff}; the value remains an inf--sup of KL-type objects
	(Remark~\ref{rem:dv}). What is lost is only the single fixed reference and the
	projection geometry it induces. The mechanism is visible at first
	order in $\theta$: expanding
	$\psi_\theta=\psi^{(0)}+\theta\psi^{(1)}+O(\theta^2)$,
	$\varphi_\theta=\varphi^{(0)}+\theta\varphi^{(1)}+O(\theta^2)$ around the standard
	bridge pair and matching the $O(\theta)$ terms of \eqref{eq:psi-rs} and
	\eqref{eq:G-fwd} gives linear equations with explicit sources: $\psi^{(1)}$
	solves the standard backward PDE with source $V\psi^{(0)}$, and $\varphi^{(1)}$
	solves the standard forward PDE with the negative source
	$-\frac{\kappa_t}{2}\|\nabla\log\psi^{(0)}\|^2\,\varphi^{(0)}$.
	
	\begin{proposition}[First-order density correction]
		\label{prop:first-order}
		The optimally controlled density expands as
		\begin{equation}
			\label{eq:p1}
			\begin{split}
				p_\theta^*(t,x)&=\varphi^{(0)}\psi^{(0)}
				+\theta\bigl[\varphi^{(0)}\psi^{(0)}\log\psi^{(0)}\\
				&\qquad+\varphi^{(0)}\psi^{(1)}+\varphi^{(1)}\psi^{(0)}\bigr]
				+O(\theta^2),
			\end{split}
		\end{equation}
		and, writing $p_0:=\varphi^{(0)}\psi^{(0)}$ and $\chi:=\psi^{(1)}/\psi^{(0)}$, the
		correction $p_1$ satisfies the linearised Fokker--Planck equation
		\begin{equation}
			\label{eq:p1-cons}
			\begin{split}
				\partial_t p_1&=-\nabla\cdot\bigl(\bigl(f+\kappa_t\nabla\log\psi^{(0)}\bigr)p_1\bigr)
				+\tfrac{\kappa_t}{2}\Delta p_1\\
				&\quad-\kappa_t\nabla\cdot\bigl(\bigl(\nabla\log\psi^{(0)}+\nabla\chi\bigr)p_0\bigr),
			\end{split}
		\end{equation}
		with $\int p_1(t,x)\,dx=0$ for all $t$.
	\end{proposition}
	
	\begin{proof}
		With
		$\nabla\log\psi_\theta=\nabla\log\psi^{(0)}+\theta\nabla\chi+O(\theta^2)$,
		matching the $O(\theta)$ terms of the exact Fokker--Planck equation gives
		\eqref{eq:p1-cons}; the integral constraint follows by differentiating
		$\int p_\theta^*\,dx=1$ at $\theta=0$.
	\end{proof}
	
	The first-order effect of risk sensitivity is therefore exactly localized:
	forward mass is depleted where the {score energy}
	$\mathcal I[\psi^{(0)}]:=\|\nabla\log\psi^{(0)}\|^2$ of the standard bridge is large (near saddles, ridges, or the ambiguous zone between modes of a mixture target), so the ensemble commits to a basin earlier and crosses
	transition zones faster: a design handle on which trajectories a generative
	sampler avoids.
	
	\section{Conclusion}
	\label{sec:conclusion}
	
	We have constructed a risk-sensitive deformation of the Schr\"odinger bridge
	that preserves exact terminal matching and the linearising Cole--Hopf
	structure, and proved that its value admits no fixed-reference KL dual: the
	obstruction is an explicit $\theta/2$ heat-equation defect, ruling out every
	endpoint-regular reference; per-$\theta$ representations are circular, and no
	single reference serves two distinct $\theta$. Open problems include a general existence theory,
	contraction of the fixed-point iteration for $\theta>0$, the limit
	$\theta\to1$, and a duality theory for $\cJ_\theta$.
	
	\section{Acknowledgments}
	The author used large language models to assist with editing, code
	refactoring, and the organization of computational experiments. All mathematical derivations, scientific claims, and final code were independently checked, validated, and approved by the author.
	
	\bibliographystyle{IEEEtran}
	\bibliography{references}
	
\end{document}


